\documentclass[10pt,a4paper,twocolumn]{article}
\usepackage[T1]{fontenc}
\usepackage[utf8]{inputenc}
\usepackage{mathptmx}
\usepackage[margin=20mm,bottom=23mm]{geometry}
\usepackage{microtype}
\usepackage{booktabs,array,tabularx}
\usepackage[font=small,labelfont=bf]{caption}
\usepackage{enumitem}
\usepackage{titlesec}
\usepackage{flushend}
\usepackage{xurl}
\usepackage[hidelinks]{hyperref}

\setlength{\columnsep}{6mm}
\setlength{\emergencystretch}{1em}
\setlength{\tabcolsep}{5pt}
\renewcommand{\arraystretch}{1.16}
\titleformat{\section}{\large\bfseries\raggedright}{\thesection}{0.65em}{}
\titleformat{\subsection}{\normalsize\bfseries\raggedright}{\thesubsection}{0.65em}{}
\titlespacing*{\section}{0pt}{11pt plus 2pt minus 1pt}{6pt}
\titlespacing*{\subsection}{0pt}{9pt plus 2pt minus 1pt}{4pt}
\setlist{nosep,leftmargin=*}
\newcolumntype{Y}{>{\raggedright\arraybackslash}X}
\newcommand{\code}[1]{\texttt{#1}}
\newcommand{\source}[1]{\path{#1}}
\hypersetup{
  pdftitle={QEMU on Zephyr: Extending the Boundaries of POSIX, Simulation and Virtualization},
  pdfauthor={Chao Liu},
  pdfsubject={QEMU porting and Zephyr POSIX compatibility assessment},
  pdfkeywords={Zephyr, POSIX, QEMU, AArch64, TCG, virtualization}
}

\title{\bfseries QEMU on Zephyr: Extending the Boundaries of\\
  POSIX, Simulation and Virtualization}
\author{Chao Liu\\[2pt]\normalsize Process Mission}
\date{}

\begin{document}
\maketitle

\begin{abstract}
Porting QEMU to Zephyr subjects an embedded operating system's POSIX layer
to the requirements of a substantial emulator: threads, synchronization,
clocks, file descriptors, image loading, executable memory and runtime
lifecycle. This paper examines that porting process and uses the resulting
workloads to assess Zephyr's practical POSIX compatibility. The assessment
separates reusable POSIX services, semantic adaptations, direct kernel
mechanisms and deliberately excluded QEMU features. It identifies concrete
boundaries in regular-file positional I/O, file metadata, memory protection,
signals and process assumptions, while distinguishing those boundaries from
GLib dependencies and the Linux guest ABI. The implementation retains QEMU
device models and its Tiny Code Generator (TCG), adds an Arm EL2 execution
backend, and executes a limited set of static Linux programs through a syscall
adapter. Recorded Linux, firmware and user-program executions demonstrate
the selected host contracts working together on an emulated Arm platform.
The resulting architecture extends Zephyr toward simulation workloads and
Type-1 hypervisor deployment on application-class SoCs. Physical-board
performance, standards-wide conformance and production isolation remain
outside the reported evidence.
\end{abstract}

\noindent\textbf{Keywords:} Zephyr; POSIX compatibility; QEMU porting;
AArch64; TCG; Type-1 hypervisor.

\section{QEMU as a POSIX compatibility workload}
\label{sec:introduction}

A large emulator exercises an operating-system interface through sustained
interactions among subsystems. CPU realization allocates and initializes
objects, image loaders combine metadata queries with positioned reads,
translation engines manage executable code and invalidation, and execution
loops coordinate timers, wakeups and shutdown. These interactions expose
assumptions that an isolated successful API call can leave unexamined.

QEMU's Unix host interface makes extensive use of POSIX facilities,
alongside its own portability abstractions and library dependencies. Its
architecture also separates complete-machine execution from Linux
user-program execution~\cite{bellard}. That division supplies several
workloads for studying the host boundary.

Zephyr provides a configurable POSIX subset within an RTOS. Applications and
kernel code are normally linked into one artifact; the enabled profile and
libc determine the services available to that application. Portability
therefore depends on configuration, API semantics, resource limits and the
runtime model in which calls occur~\cite{zephyr-posix}.

The project uses QEMU as an integration workload with three objectives:
\begin{enumerate}
  \item Identify which host requirements can use enabled Zephyr POSIX services.
  \item Make semantic adaptations and unsupported requirements explicit.
  \item Demonstrate the retained emulator components through Linux, firmware
    and Linux ABI program execution.
\end{enumerate}

The assessment is specific to the pinned source, selected configuration and
exercised paths. A successful guest boot supplies application-level evidence
for the configured system. Standards-wide conformance requires a separate
requirement inventory and semantic test suite.

\section{Assessment frame}
\label{sec:frame}

The study distinguishes three interfaces: the host POSIX API, QEMU and
library contracts, and the guest Linux ABI. For each requirement, it records
the pinned implementation, the selected adaptation, the evidence that
exercises it, and the limits of that evidence. Source inspection establishes
details such as an explicit \code{ENOSYS} return. Runtime tests establish
behavior along their actual execution paths.

\subsection{The host POSIX interface}

This is the interface QEMU itself consumes inside Zephyr: pthread operations,
clocks, file operations and related host services. Its assessment concerns
the behavior that Zephyr and the selected libc supply to QEMU. The application
enables POSIX threads and files, compiler TLS and Picolibc, and configures
finite pools for threads, mutexes, condition variables and descriptors.
The artifact records these selections in
\source{apps/qemu_linux/prj.conf} and \source{zephyr/Kconfig}.

\subsection{QEMU and library contracts}

The QEMU Object Model (QOM), MemoryRegion, the big QEMU lock (BQL),
read-copy-update (RCU), GLib containers and TCG are QEMU or library contracts.
Their host requirements are assessed separately from POSIX API coverage.
The port's GLib subset has differential tests. Disabling a QEMU subsystem
such as migration describes the selected port surface; additional evidence
would be needed to attribute that exclusion to a missing Zephyr POSIX service.

\subsection{The guest Linux ABI}

In user mode, a guest program issues Linux AArch64 syscall numbers. A project
adapter interprets those requests and invokes available host operations.
Zephyr has its own build-generated kernel syscall identifiers and a POSIX
library interface~\cite{zephyr-implementation}. The Linux ABI adapter is an
additional compatibility layer, and its syscall tests evaluate that layer.
Its implementation is \source{src/qemu/ports/zephyr/user-syscall.c}.

\section{Porting method}
\label{sec:porting}

\subsection{Establish a reproducible build boundary}

The build starts from pinned QEMU, Zephyr, dtc and zlib revisions. It exports
those sources into \source{build/sources/}, applies ordered patches, and
overlays new implementation files. Zephyr CMake selects the QEMU sources and
invokes the original QAPI, trace and instruction-decoder generators. The
upstream checkouts remain clean. The manifest is \source{west/west.yml}; the
module build is \source{zephyr/CMakeLists.txt}.

The port defines a Zephyr host environment in QEMU's platform abstraction.
Header selection, page-size discovery and macro names require explicit
handling: Zephyr's numeric configuration macros and QEMU's macro conventions
must retain their respective meanings. Aligned allocations use the selected
libc with size rounding and overflow checks. These are compiler, build and
C-library adaptations adjacent to the POSIX interface. They are implemented
in the host configuration and the host-header and allocation patches.

\subsection{Retain an executable QEMU component surface}

The system source set retains ARM CPUs, QOM/qdev, MemoryRegion dispatch,
PL011, a software GICv3 and the ARM loaders. A device-model probe exercises
QOM construction and PL011 register accesses on Zephyr. This supplies a
runtime checkpoint for the object, allocation and device-dispatch layers
before interpreting the results of a complete guest workload. The probe is
\source{apps/qemu_probe/}; the machine implementation is
\source{src/qemu/ports/zephyr/machine.c}.

The selected GLib interface supplies data structures and utility operations
needed by that source set. Its differential test executes the local
implementation and the development host's GLib with matching inputs,
including allocation-overflow behavior. The result supports those exercised
library semantics. Full GLib compatibility and POSIX conformance would each
require their own evaluations. The implementation and tests are recorded in
\source{src/qemu/ports/zephyr/glib/} and \source{tests/glib/}.

\subsection{Give the embedded runtime an owner}

The application uses \code{pthread\_create()} to start a single POSIX worker
with explicit thread attributes. QEMU mutex and condition-variable wrappers call
Zephyr's pthread functions. The worker owns model operations, CPU execution,
timers and deferred reclamation. The port preserves BQL ownership across
execution transitions and reclaims deferred RCU callbacks at quiescent points
outside read-side sections. These rules are implemented by the application
and \source{src/qemu/ports/zephyr/os.c}.

This ownership policy is also a semantic precondition for several adapters.
Private file descriptors permit a seek/read/restore implementation of
positioned I/O. Single-owner initialization permits bounded handling of
one-time initialization and RCU readers. Those policies must be revisited
when adding QEMU threads or shared-descriptor concurrency.

\subsection{Connect image loading to real storage}

Outer QEMU supplies VirtIO Block storage. Zephyr mounts Ext2 read-only at
\source{/images}, and the loader reaches it through the file adapter. Kernel,
initramfs and ELF data use real filesystem operations. GLib mapped-file
objects are represented by private file snapshots whose references and
lifetime are managed by \source{src/qemu/ports/zephyr/file.c}.

The Ext2 adaptation makes inode synchronization conditional on mount
writability. This is a filesystem integration requirement beneath POSIX:
loader close and synchronization paths must respect the mounted medium's
read-only contract. The implementation is preserved as an ordered Zephyr
patch in the artifact.

\subsection{Add execution and an explicit command lifecycle}

The system runtime creates the machine, selects its compiled accelerator,
loads the image and services execution returns. The user runtime uses a
separate source selection with upstream user TCG and Linux ELF loading.
Both are exposed through Zephyr shell commands that copy arguments before
waking the worker.

The shell owns UART input and uses a bypass callback while a guest runs.
Ctrl-] requests termination, and completion restores the prompt. Automatic
startup calls the same handler through \code{shell\_execute\_cmd()}~\cite{shell}.
A worker-local exit target converts QEMU loader exits into command return
status, giving embedded execution an explicit lifecycle within the surrounding
RTOS. The console adapter is \source{src/qemu/ports/zephyr/console-uart.c}.

\section{Zephyr POSIX compatibility assessment}
\label{sec:assessment}

\subsection{Findings by host requirement}

Table~\ref{tab:posix} associates host requirements with the pinned interface,
the port treatment and the relevant evidence. ``Adapted'' identifies behavior
implemented by the port; ``kernel mechanism'' identifies a direct Zephyr
facility used at that boundary. A successful execution establishes the
behavior of the combined system. Claims about unadapted Zephyr POSIX behavior
require evidence for that interface itself.

\begin{table*}[t]
\caption{Workload-specific POSIX assessment at the pinned Zephyr revision.}
\label{tab:posix}
\small
\begin{tabularx}{\textwidth}{@{}>{\raggedright\arraybackslash}p{0.20\textwidth}YY@{}}
\toprule
Host requirement & Pinned interface or configuration & Port treatment and evidence \\
\midrule
Worker and synchronization & Pthreads enabled; finite thread, mutex and
condition-variable pools & Real pthread calls support the owner and used
locking paths; guest execution supplies integration evidence. \\
Clock access and waiting & POSIX clock calls and Zephyr sleep/timer facilities &
Clock queries are reused; deadline conversion and wakeups are adapted;
clock and scheduling checks exercise selected paths. \\
Basic file operations & POSIX filesystem APIs over Zephyr VFS &
Open, read, seek, status and close load real images and user ELF files. \\
Positioned file I/O & Offset dispatch accepts shared-memory descriptors &
Private-descriptor seek/I/O/restore supplies positioned regular-file access. \\
File permissions & Generic filesystem status reports type and size without
execute bits & Loader policy uses available metadata and retains ELF and
mapping validation. \\
Mapping and protection & Host \code{mprotect()} returns \code{ENOSYS}; POSIX
mapped files disabled in user configuration & Kernel memory APIs provide JIT
aliases; QEMU metadata and helpers enforce guest access rules. \\
Processes and signals & Linked application model; several process-signal APIs
return \code{ENOSYS} & Worker-local exit handling and synthetic guest
termination implement the selected lifecycle. \\
General QEMU facilities & Capabilities depend on selected QEMU sources &
Migration, hotplug, general coroutines and multiple vCPUs are excluded by
the port's implemented profile. \\
\bottomrule
\end{tabularx}
\end{table*}

\subsection{Pthreads and time}

Worker creation uses explicit pthread attributes to select \code{SCHED\_RR}
and supply the thread's memory region. Mutex and condition operations call
the corresponding pthread functions. Timed condition waits use an absolute
deadline: the wrapper reads \code{CLOCK\_REALTIME} through
\code{clock\_gettime()}, adds the requested duration and normalizes
nanoseconds.

These bindings invoke the host APIs and check their return codes. The
execution record demonstrates creation of the owner, model locking and
progress of an independent Zephyr thread while a guest executes. Pthread
cancellation, priority interactions, multi-owner contention and every
timed-wait path have not been exhaustively exercised. Source binding and
observed behavior are therefore reported at different levels of granularity.

Clock and scheduler integration also uses direct kernel services. Exit
requests originate from Zephyr timers; halted execution uses wait/kick
synchronization. These mechanisms allow the selected QEMU loop to cooperate
with Zephyr's scheduler~\cite{scheduling}. An observed wakeup or advancing
timer establishes progress in the tested run. Worst-case latency requires
separate analysis.

\subsection{Positioned I/O and descriptor semantics}
\label{sec:positioned-io}

At the pinned revision, Zephyr's \code{pread()} and \code{pwrite()} pass an
explicit offset into \code{zvfs\_rw()}. The descriptor layer's
\code{supports\_pread\_pwrite()} accepts shared-memory descriptors. A positioned
operation on a regular file therefore reaches an \code{ENOTSUP}
result~\cite{device-io,zvfs}. QEMU's ELF loading path requires positioned
regular-file reads, making descriptor semantics material to the port even
though the API symbols are present.

The project's adapter saves the current offset, seeks, performs I/O and
restores the offset. Its contract relies on QEMU owning these descriptors on
one worker. ELF loading and private file-mapping checks exercise this path.
Concurrent positioned access through a shared open-file description would
require a stronger implementation and separate tests. The evaluation credits
the port adapter for the supported behavior.

\subsection{File metadata and loader policy}

The pinned \code{stat()} path obtains a Zephyr filesystem entry and reports
its type and size. Its mode contains \code{S\_IFREG} or \code{S\_IFDIR},
without execution permission bits~\cite{zephyr-stat}. The selected Linux ELF
loader cannot apply its conventional execute-bit policy using that metadata.
The adaptation accepts an explicitly requested regular executable and retains
ELF architecture and address-range validation.

This finding concerns the generic filesystem/POSIX boundary used by the
port. On-disk permission metadata and the information exposed by a host API
are separate concerns. Consumers need the relevant policy information and
semantics to be carried through the interface they actually use.

\subsection{Memory protection and fault delivery}

The host protection function \code{mprotect()} returns \code{ENOSYS} at the
pinned revision~\cite{zephyr-mprotect}. The user configuration explicitly
disables POSIX mapped files. The memory behavior demonstrated by this port
is implemented through a defined combination of kernel mechanisms and QEMU
metadata.

For generated code, an aligned backing buffer receives distinct writable
and executable mappings using \code{k\_mem\_map\_phys\_bare()}. The TCG
adapter uses the upstream cache-synchronization machinery to publish code
through these aliases. The configured buffer size is 16~MiB. The system
adapter is \source{src/qemu/ports/zephyr/tcg.c}.

For Linux user memory, the runtime supplies a bounded 64~MiB region and
implements mapping operations with QEMU page metadata. Generated user loads
and stores pass through checked helpers; the integration includes atomic
and bulk-memory paths such as DC ZVA. Writes to translated code invoke QEMU's
invalidation mechanism. Invalid accesses terminate the emulated process with
a signal-derived status. The memory adapter is
\source{src/qemu/ports/zephyr/user-memory.c}.

The resulting fault tests validate the project adapter. A native POSIX
\code{mprotect()} assessment would instead need to exercise that Zephyr
entry point and its prescribed semantics directly.

\subsection{Signals, identity and process lifetime}

The pinned signal implementation provides signal-set operations while
several process-signal functions, including \code{sigaction()} and
\code{kill()}, return \code{ENOSYS}~\cite{zephyr-signals}. Enabling
signal-related configuration alone therefore cannot establish the
host-fault-delivery model expected by a Linux user emulator. The port maps
the selected no-mask-save jump paths to \code{setjmp}/\code{longjmp} and
detects guest faults through its memory helpers.

Guest Linux PID, UID and signal contracts are distinct from host pthreads
and Zephyr task identity. The user adapter assigns a positive process ID per
launch and a virtual root identity. Unsupported Linux calls return
\code{ENOSYS}. The implemented guest ABI excludes installed signal handlers
and guest threads. Zephyr separately offers pthread and socket interfaces,
so their omission from this guest ABI is a port-scope decision.

System execution retains its initialized machine and global QEMU state after
the loop returns. A new VM requires a Zephyr reboot. User execution closes
descriptors and resets mappings, CPU and translation state for successive
programs. These rules give the embedded runtime explicit ownership of the
lifecycle normally supplied by a standalone emulator's host process.

\subsection{Scope of the assessment}

Threads, selected synchronization paths, clocks and basic file operations
form a usable foundation for the retained workload. Positioned I/O, metadata,
memory protection, signals and process assumptions identify interfaces that
need semantic qualification or adaptation in this build. Finite resource
settings and single-owner execution are part of the result.

A broader maturity assessment would require a declared POSIX profile, a
complete requirement inventory, targeted semantic tests and concurrency
coverage. The current evidence supports workload compatibility for the
documented configuration and identifies concrete requirements for API
qualification and implementation.

\section{Simulation and virtualization capabilities}
\label{sec:capabilities}

\subsection{QEMU machine and TCG execution}

QEMU's TCG translates guest instructions into host code and reuses translation
blocks associated with CPU state. System emulation uses a software MMU and
routes MMIO through device models~\cite{tcg}.

The project combines that engine with ARM CPU objects, PL011, a software
GICv3 and the ARM Linux/ELF loaders. Its \code{zephyr-virt} machine provides
one vCPU and 256~MiB RAM at \code{0x40000000}, with PL011 at
\code{0x09000000}. System TCG runs on an EL1 Zephyr host. A one-millisecond
timer requests execution exits so the owner can service timers and console
work. This supplies a software simulation path within the RTOS.

\subsection{EL2 execution and the Type-1 direction}

Arm virtualization provides EL2 controls for trapping and Stage-2 address
translation. These mechanisms let a hypervisor define a guest's resource
view and handle accesses that require emulation~\cite{arm-virtualization}.

The native profile places Zephyr at non-VHE EL2 and connects the QEMU
\code{zephyr} accelerator to a kernel executor named \code{zhv}. The executor
handles guest entry/exit, Stage-2 mappings, registers, TLS, SIMD/FP state and
timers. The QEMU ARM adapter dispatches MMIO, system-register and PSCI exits
and uses software GICv3 state for interrupt delivery. Native CPU realization
checks the relevant host MIDR fields against the selected model.

This architecture gives Zephyr the execution control and device-model layer
needed to pursue a Type-1 hypervisor role on application-class SoCs. Such a
deployment places Zephyr in control of the physical virtualization interface.
The present experiments use an outer QEMU \code{virt} platform with
TCG~\cite{virt}, so the native backend exercises virtualized Arm mechanisms
within that development environment. A physical deployment additionally
requires board support, memory and interrupt configuration, and hardware
validation.

The executor source is under
\source{src/zephyr/arch/arm64/core/hypervisor/}; the QEMU ARM adapter is
\source{src/qemu/target/arm/zephyr.c}. System accelerator selection and user
mode are separate firmware builds.

\subsection{Linux user execution as an additional workload}

The user path reuses QEMU user TCG and its Linux ELF loader. At \code{svc},
it reads the syscall number from \code{x8}, supplies arguments from
\code{x0} through \code{x5} to the adapter and writes the result to \code{x0}.
Files, private memory mappings, clocks, entropy and selected process metadata
are supported. QEMU's user-emulation documentation supplies the general
model~\cite{qemu-user}; the Zephyr-specific syscall surface is defined by
the project implementation.

The default program is a static Linux ABI ELF built with the installed SDK.
A larger regression program exercises file I/O, mappings and faults. These
programs extend the workload beyond a kernel boot, while retaining the
distinction between host POSIX compatibility and guest Linux ABI support.

\section{Evaluation method and recorded results}
\label{sec:evaluation}

\subsection{Source and environment}

The implementation baseline is repository commit
\source{2b1161806afe14fe57dec97234f37aac1ea84e5e}. The evaluation combines source
inspection with recorded executions and checked-in acceptance workloads.
Run reports are versioned in \source{docs/validation.md}. Raw console logs
are generated locally and excluded from Git. The record establishes successful
finite executions; repetition counts for failure-rate or timing-distribution
estimates are not available.

Table~\ref{tab:environment} records the principal inputs. The manifest also
pins dtc and zlib, and \source{docs/guest-assets.md} records image hashes.
Default configured capacities include 512~MiB outer RAM, 256~MiB system guest
RAM, a 64~MiB user address region, a 16~MiB TCG code buffer, and a 32~MiB native
or 96~MiB TCG/user C-library arena. These are configuration values; peak
memory consumption requires separate measurement.

\begin{table*}[t]
\caption{Recorded source, environment and workload inputs.}
\label{tab:environment}
\small
\begin{tabularx}{\textwidth}{@{}>{\raggedright\arraybackslash}p{0.22\textwidth}Y@{}}
\toprule
Input & Recorded configuration \\
\midrule
QEMU source & \source{1df256f5968e9f7c3c4533a1383b071c044a36d6}; version 11.1.50 development tree \\
Zephyr source & \source{ba25413e5b6b2661a71db24888f6b89274f1480d}; version 4.4.99 development tree \\
Toolchain & Zephyr SDK 1.0.1, AArch64 GNU \\
Outer emulator & SDK QEMU 10.0.2; records also include system QEMU 10.2.2 \\
Development hosts & macOS Apple Silicon and Ubuntu 24.04 AArch64 container \\
System workloads & Linux 6.4.16 Image, pinned BusyBox initramfs, ELF and raw firmware \\
User workloads & Default hello, Linux ABI regression program and static glibc program \\
\bottomrule
\end{tabularx}
\end{table*}

\subsection{Evidence and attribution}

Table~\ref{tab:evidence} associates the observed behavior with the interface
being evaluated. The Linux checker reads the architectural timer interrupt
count, sleeps in the guest and confirms that the count advances. It runs a
busy loop for approximately eight guest-uptime seconds and queries the
restored Zephyr shell after poweroff. Positive EL0 and MMU-on samples show
that execution reached those states. The counters sample execution returns;
they do not measure instruction totals or processor utilization. The checker
is \source{apps/qemu_linux/check.py}.

\begin{table*}[t]
\caption{Recorded functional evidence and its attribution.}
\label{tab:evidence}
\small
\begin{tabularx}{\textwidth}{@{}>{\raggedright\arraybackslash}p{0.20\textwidth}YY@{}}
\toprule
Experiment & Observed behavior & Compatibility evidence and limit \\
\midrule
Linux: both backends, A53/A57/A72 & Console I/O, timer IRQ growth, EL0/MMU
samples, host progress and shutdown & Integrated host interface works for
recorded profiles; individual POSIX APIs require separate characterization. \\
ELF/raw firmware: both backends & Segments, initialized data, BSS, EL1 entry,
PSCI shutdown and reboot & Loader, file and execution paths operate together. \\
User ABI: A53/A57/A72 & Arguments, environment, files, errno, mappings, code
modification, time and entropy & Exercises the combined user runtime and adapters. \\
User fault cases & Invalid addresses, read-only writes, DC ZVA protection,
unmapped access and relaunch & Exercises software guest-memory enforcement
and synthetic termination. \\
Default user Make workflow & Real hello ELF, arguments, environment, tracing
and clean outer exit & Confirms the advertised program is present and runnable. \\
Static glibc program & Initialization, stdio reads, allocation, clocks and
repeated output & Establishes one additional workload with specific syscall usage. \\
Filesystem and GLib checks & File-operation cases and differential library
behavior & Evidence applies to the exercised adapters and library semantics. \\
\bottomrule
\end{tabularx}
\end{table*}

The independent observer sleeps for one second between updates. The checker
requires progress and a maximum recorded gap below its 4000~ms acceptance
threshold. This is a finite-run health criterion. Worst-case integration
timing would also require analysis of initialization, interrupt masking,
allocation, device operations and the outer emulator's scheduling.

User acceptance executes real AArch64 instructions and Linux syscall numbers.
It checks private file mappings, self-modifying code, memory permissions and
termination, then starts another program. The default hello ELF also executed
directly in Ubuntu AArch64. Component evidence includes 196 payload-filesystem
checks, 184 matching GLib differential output lines and the recorded
architecture/FPU/executor cases. These counts describe the existing checks;
they provide no coverage percentage. User acceptance is implemented by
\source{tests/user/check.py} and \source{tests/user/check_default.py}.

\subsection{File updates and descriptor lifetime}

External files are assembled into a read-only Ext2 image using source-content
fingerprints and atomic replacement of the completed disk. Outer QEMU keeps
its attached disk for the session. Restarting Zephyr resets inner runtime
state, while the outer file attachment remains. The disk builder is
\source{scripts/guest_disk.py}.

An update experiment published two checked-in workloads under one path:
hello first, then the Linux ABI regression program. After the host binary and
disk were replaced, the existing session still executed hello, including
after a Zephyr reboot. A new outer QEMU process executed the ABI workload
and returned its deliberate success-test status of 7. This confirms the
documented activation boundary for that workflow and evaluates the storage
and lifecycle integration.

\section{Implications for POSIX maturity and deployment}
\label{sec:implications}

The adaptation matrix identifies contracts that determine application
portability. Regular-file positioned I/O needs appropriate offset semantics;
filesystem metadata must expose the policy information consumers require;
protection and signal interfaces need explicit qualification; and application
lifetime must be reconciled with an embedded runtime.

Future compatibility evaluations should attribute each success to its actual
provider: Zephyr POSIX, the selected libc, a kernel API or a project adapter.
They should exercise the relevant error and concurrency semantics as well as
the normal path. A QEMU feature excluded by configuration should remain
separate from an unavailable POSIX function in the resulting assessment.

The current artifact restricts system execution to one VM and one vCPU per
Zephyr boot. User execution permits one static, single-threaded program at a
time. Migration, hotplug, guest EL2/EL3, general signal delivery, dynamic
linking and guest process creation are outside the implemented profile.
The input disk is a snapshot; updated files become active after regeneration
and an outer QEMU restart.

The experiments use an emulated Arm platform. The evaluation record contains
no physical-board throughput, interrupt-latency, energy or memory-bandwidth
measurements. Security testing covers selected faults and mechanisms; input
parsing and model code still run within the Zephyr application environment.
A production isolation claim or a real-time bound requires a separate
assessment. Deployment on a high-performance SoC also requires its BSP,
memory layout, GIC, timer and exception-level configuration.

\section{Artifact and reproduction}
\label{sec:artifact}

The artifact includes pinned upstream inputs, ordered source patches, new
host and executor code, sample programs and acceptance scripts~\cite{project}.
Supported development hosts are Linux x86\_64/AArch64 and macOS Apple Silicon,
with Python 3.12 or newer. After checking out the implementation revision in
Section~\ref{sec:evaluation}, the main paths can be reproduced as follows:

\begin{quote}
\begin{minipage}{\linewidth}
\small\ttfamily
bash scripts/install-host-deps.sh\\
make setup\\
make doctor\\
make check ACCEL=zephyr CPU=cortex-a53\\
make check QEMU\_SHELL=1 ACCEL=tcg\\
make check-firmware\\
make check-firmware ACCEL=tcg\\
make check-user\\
make probe test-payload test-glib\\
make test-tools\\
make test-arch
\end{minipage}
\end{quote}

The documented A57 and A72 selections reproduce the additional CPU-model
coverage. Each run should record its repository revision, SDK and outer QEMU
versions, configuration and console transcript. The Python environment
includes dependencies from the pinned Zephyr requirements; the complete
transitive environment is not locked.

English and Chinese usage guidelines in \source{docs/guidelines.md} and
\source{docs/guidelines.zh-CN.md} specify installation, program updates and
restart procedures. The principal assessment sources are the host runtime,
file adapter, host header, module source lists and the pinned Zephyr
implementations cited in Section~\ref{sec:assessment}. Source references and
validation records were consulted on 9 October 2026.

\begin{thebibliography}{16}
\small
\setlength{\itemsep}{3pt}

\bibitem{bellard}
F. Bellard, ``QEMU, a Fast and Portable Dynamic Translator,'' in
\emph{FREENIX Track: 2005 USENIX Annual Technical Conference},
USENIX Association, 2005, pp. 41--46.
\href{https://www.usenix.org/conference/2005-usenix-annual-technical-conference/qemu-fast-and-portable-dynamic-translator}{Publication record}.

\bibitem{zephyr-posix}
Zephyr Project,
\href{https://docs.zephyrproject.org/latest/services/portability/posix/overview/index.html}{``POSIX: Overview,''}
\emph{Zephyr Project Documentation}.

\bibitem{zephyr-implementation}
Zephyr Project,
\href{https://docs.zephyrproject.org/latest/services/portability/posix/implementation/index.html}{``POSIX: Implementation Details,''}
\emph{Zephyr Project Documentation}.

\bibitem{shell}
Zephyr Project,
\href{https://docs.zephyrproject.org/latest/services/shell/index.html}{``Shell,''}
\emph{Zephyr Project Documentation}.

\bibitem{scheduling}
Zephyr Project,
\href{https://docs.zephyrproject.org/latest/kernel/services/scheduling/index.html}{``Scheduling,''}
\emph{Zephyr Project Documentation}.

\bibitem{device-io}
Zephyr Project, ``POSIX device I/O implementation,''
\source{subsys/portability/posix/options/device_io.c},
revision \source{ba25413e5b6b2661a71db24888f6b89274f1480d}.
\href{https://github.com/zephyrproject-rtos/zephyr/blob/ba25413e5b6b2661a71db24888f6b89274f1480d/subsys/portability/posix/options/device_io.c}{Pinned source}.

\bibitem{zvfs}
Zephyr Project, ``ZVFS descriptor dispatch,''
\source{lib/os/zvfs/zvfs_fdtable.c}, same revision as~\cite{device-io}.
\href{https://github.com/zephyrproject-rtos/zephyr/blob/ba25413e5b6b2661a71db24888f6b89274f1480d/lib/os/zvfs/zvfs_fdtable.c}{Pinned source}.

\bibitem{zephyr-stat}
Zephyr Project, ``POSIX filesystem implementation,''
\source{subsys/portability/posix/options/fs.c}, same revision as~\cite{device-io}.
\href{https://github.com/zephyrproject-rtos/zephyr/blob/ba25413e5b6b2661a71db24888f6b89274f1480d/subsys/portability/posix/options/fs.c}{Pinned source}.

\bibitem{zephyr-mprotect}
Zephyr Project, ``POSIX memory protection implementation,''
\source{subsys/portability/posix/options/mprotect.c}, same revision as~\cite{device-io}.
\href{https://github.com/zephyrproject-rtos/zephyr/blob/ba25413e5b6b2661a71db24888f6b89274f1480d/subsys/portability/posix/options/mprotect.c}{Pinned source}.

\bibitem{zephyr-signals}
Zephyr Project, ``POSIX signal implementation,''
\source{subsys/portability/posix/options/signal.c}, same revision as~\cite{device-io}.
\href{https://github.com/zephyrproject-rtos/zephyr/blob/ba25413e5b6b2661a71db24888f6b89274f1480d/subsys/portability/posix/options/signal.c}{Pinned source}.

\bibitem{tcg}
QEMU Project,
\href{https://www.qemu.org/docs/master/devel/tcg.html}{``Translator Internals,''}
\emph{QEMU Documentation}, version 11.1.50.

\bibitem{arm-virtualization}
Arm Limited, \emph{Armv8-A virtualization}, document 102142, Issue 01, 2019.
\href{https://developer.arm.com/-/media/Arm\%20Developer\%20Community/PDF/Learn\%20the\%20Architecture/Armv8-A\%20virtualization.pdf?revision=a765a7df-1a00-434d-b241-357bfda2dd31}{Official document}.

\bibitem{virt}
QEMU Project,
\href{https://www.qemu.org/docs/master/system/arm/virt.html}{```virt' generic virtual platform,''}
\emph{QEMU Documentation}.

\bibitem{qemu-user}
QEMU Project,
\href{https://www.qemu.org/docs/master/user/main.html}{``QEMU User space emulator,''}
\emph{QEMU Documentation}.

\bibitem{project}
\href{https://github.com/processmission/qemu-on-zephyr}{QEMU on Zephyr},
source and validation artifact.
Implementation baseline: \source{2b1161806afe14fe57dec97234f37aac1ea84e5e}.

\end{thebibliography}
\end{document}
